\subsection{The missing arithmetic inputs at exact class-four constants}

\begin{result}{E.80}{Proved}{The exact four-place weight-three matching map}
The continuous restriction map
\[
\begin{split}
 H^1(G_{\mathbb Q},\widehat{\mathbb Z}(3))\longrightarrow{}&
 H^1(G_{\mathbb Q_2},\mathbb Z_7(3))
 \oplus H^1(G_{\mathbb Q_3},\mathbb Z_{13}(3))\\
 &\oplus H^1(G_{\mathbb Q_3},\mathbb Z_2(3))
 \oplus H^1(G_{\mathbb R},\mathbb Z_2(3))
\end{split}
\tag{E.80.1}
\]
is surjective. The four targets have orders $7,13,2,2$.
This independently confirms SECTIONS\_OVER\_Q A.60.
Consequently, above the distinguished integral class-two
section of an ordinary pair whose valuations at two and three
are $(0,1)$ and $(1,0)$, respectively, any global integral
class-three lift can be changed to match the ordinary local
point classes at these four indicated places and coefficients.
Both central weight-three coordinates can be prescribed
independently. No local section theorem is used.
\end{result}
\begin{proof}
For a finite place $q\ne\ell$, the integral tame calculation is
\[
 H^1(G_{\mathbb Q_q},\mathbb Z_\ell(3))
       =\mathbb Z_\ell/(q^3-1)\mathbb Z_\ell.
\]
The inertia invariant term is zero because multiplication by
$q^2-1$ is injective on $\mathbb Z_\ell$. At the real place,
complex conjugation acts by minus one, so the first cohomology
of its order-two group is $\mathbb Z_2/2$. This proves the
orders of the displayed groups.

Let $S_\ell=\{\ell,\infty\}$. A.22 proves
$H^2(G_{\mathbb Q,S_\ell},\mathbb Z_\ell(3))=0$,
including $\ell=2$. The actual coefficient sequence therefore
makes
\[
 H^1(G_{\mathbb Q,S_\ell},\mathbb Z_\ell(3))
 \longrightarrow H^1(G_{\mathbb Q,S_\ell},\mathbb F_\ell(3))
 \tag{E.80.2}
\]
surjective. We produce finite classes with the required local
values and lift them through this map; we do not infer integral
lifting merely from a finite local calculation.

For $\ell=7,13$, put $K=\mathbb Q(\zeta_\ell)$ and define the
actual cyclotomic $S_\ell$-unit
\[
 u_\ell=\prod_{a=1}^{\ell-1}(1-\zeta_\ell^a)^{a^2}.
 \tag{E.80.3}
\]
For $\sigma_b(\zeta_\ell)=\zeta_\ell^b$, reindexing the
product gives
$[\sigma_bu_\ell]=b^{-2}[u_\ell]$ in
$K^\times/(K^\times)^\ell$. Every factor is a unit away from
$\ell$. Its Kummer torsor is consequently unramified there.
Over $K$ identify the coefficient module with
$\mathbb F_\ell(3)=\mu_\ell\otimes\mathbb F_\ell(2)$.
The additional character $b^2$ cancels the inverse-square
character of (E.80.3). The resulting cohomology class is
$\operatorname{Gal}(K/\mathbb Q)$-invariant.

Restriction in the finite Hochschild--Serre sequence is an
isomorphism onto these invariants: the quotient group has order
$\ell-1$, invertible on the coefficient module, and averaging
kills its positive-degree cohomology. We therefore obtain a
class in the right side of (E.80.2). This also specifies its
descent, not just a class over the cyclotomic field. The Kummer
identification used here is the actual radical-torsor map of
\citet[Theorems 5.30 and 5.32, pp.75--76]{RootMilneFT}.

At $(q,\ell)=(2,7)$ or $(3,13)$, $q$ has order three modulo
$\ell$. A completion of $K$ at $q$ is therefore the unramified
degree-three extension of $\mathbb Q_q$. Write $\tau$ for the
residue of the polynomial variable in the following exact
finite fields, and take $\overline\zeta_\ell=\tau^2$:
\[
\begin{array}{c|c|c|c}
 (q,\ell)&\text{residue field}&\overline u_\ell&
       \overline u_\ell^{(q^3-1)/\ell}\\ \hline
 (2,7)&\mathbb F_2[\tau]/(\tau^3+\tau+1)&\tau^2&\tau^2\\
 (3,13)&\mathbb F_3[\tau]/(\tau^3-\tau-1)&
                 2\tau+2\tau^2&2+2\tau^2
\end{array}
\tag{E.80.4}
\]
Both cubics have no root in their base field, hence are
irreducible. In both cases $\tau^2$ has the claimed prime order,
and the last entry is not one. All multiplications in this table,
as well as the complete inverse-square eigencharacter identities,
were executed in \path{code/check_weight3_matching.py}; its
output gives the literal coefficient triples. The cyclic
multiplicative group of the residue field then shows that the
unit is not an $\ell$th power. Hensel lifting identifies this
with a nonzero unramified local Kummer class. Restriction to the
degree-three extension is injective on the base local class,
since restriction followed by corestriction multiplies by three.
The integral lift from (E.80.2) thus maps nontrivially to the
order-$\ell$ target, and generates it.

For $\ell=2$ the finite coefficient module is trivial. The
two square classes $-1,2$ define classes unramified outside
$\{2,\infty\}$, and both lift by (E.80.2). At $q=3$ their
unit residues are both nonsquares, so both give the nonzero
unramified quadratic class. At the real place, only $-1$ is
negative. Coefficient reduction is injective on each of the
order-two integral local first-cohomology groups. With columns
$(-1,2)$ and rows $(\mathbb Q_3,\mathbb R)$, the actual matrix is
\[
                     \begin{pmatrix}1&1\\1&0\end{pmatrix}
                   \quad\text{over }\mathbb F_2.
 \tag{E.80.5}
\]
Its determinant is one. Thus the two targets at coefficient two
can be prescribed independently. Continuous cochains into a
product of the Tate modules are precisely products of continuous
cochains. The same holds for cocycles and coboundaries, so first
cohomology with $\widehat{\mathbb Z}(3)$ is the product of its
prime components. Combining the three coefficient-prime choices
proves (E.80.1).

For the section consequence, the specified local class-two
restrictions already match actual point sections. At $(2,7)$
and $(3,13)$ the local central weight-two groups are zero. At
$(3,2)$ the distinguished class has been selected by B.101;
the nonzero valuation is one, so its zero third obstruction is
exactly the strong class-two matching condition of B.24. At the
real place the ordinary abelian signs select an ordinary
component, and the weight-two first cohomology is zero.
Choose local conjugators aligning these lower sections. The
differences from their actual point lifts at class three are
classes in the two copies of the four target groups in
(E.80.1). Surjectivity supplies a global central twist with
their negatives. After this twist all four local section
classes match. At the real place a possible lower stabilizer
changes the chosen discrepancy representative, but not this
existence argument: we correct any one aligned representative.
\end{proof}

\begin{result}{E.81}{Proved}{The original exact class-four valuation bound}
Suppose an ordinary completed pair has valuation patterns $(0,1)$
at two and $(1,0)$ at three, and at every prime $q>3$ either
$(0,0)$ or one of
\[
 (M_q,0),\quad(0,M_q),\quad(-M_q,-M_q),\qquad
 M_q=24e_q(q^2-1)(q^3-1),
 \tag{E.81.1}
\]
where the $e_q$ are nonnegative integers tending to zero in
$\widehat{\mathbb Z}$. Its distinguished integral class-two section
has a lift to the actual integral class-four arithmetic quotient.
This independently confirms the precise constant and assertion of
SECTIONS\_OVER\_Q A.61. It does not assert ordinary local matching
of this class-four lift at the new support primes.
\end{result}
\begin{proof}
E.65 applies with $e_q(q^3-1)$ in place of its $e_q$, since this
sequence also tends to zero in each $\mathbb Z_\ell$. It gives an
integral class-three section. Apply E.80 in its two central
coordinates to make this section match the actual point sections
at $(q,\ell)=(2,7),(3,13),(3,2)$ and at the real place for
coefficient two. We prove that its class-four obstruction vanishes.

Fix $\ell$ and a fresh support prime $q>3$, $q\ne\ell$.
Write $m=M_q$ for positive support and $m=-M_q$ for negative
support. Use the degree-graded free rational Lie algebra on
$X=\log x,Y=\log y$, truncated after degree four. Let $L_i$ be
its integral free Lie lattice in degree $i$. Choose arithmetic
Frobenius $F$ and the tame $\ell$-generator $T$, and write $U,V$
for the geometric coordinates of their images under the class-three
section. The actual tangential action satisfies
\[
 \sigma(x)=x^q,\qquad \sigma(y)=f^{-1}y^qf,
       \quad f\in\Gamma_2,
 \qquad U\sigma(V)U^{-1}=V^q.
 \tag{E.81.2}
\]
The action is unramified at $q$, and the equality is at class three.
There is no replacement of the actual path element $f$ by its graded
action.

We record explicitly the denominator estimates used below. Put
$C=[x,y]$, $D=[x,C]$, $E=[y,C]$; reversing the sign of the
weight-three basis from B.23 changes none of these lattices.
Expansion of the finite series
$\exp Z=\sum_{j=0}^4Z^j/j!$ and
$\log(1+Z)=\sum_{j=1}^4(-1)^{j+1}Z^j/j$ gives
\[
\begin{array}{c|ccc}
 &\text{degree two}&\text{degree three}&\text{degree four}\\\hline
 \log C&L_2&\frac12L_3&\frac1{12}L_4\\
 \log D,\log E&0&L_3&\frac12L_4
\end{array}
\tag{E.81.3}
\]
and for arbitrary integral Hall coordinates
\[
 Z=\log(y^b x^aC^hD^dE^e):\qquad
       Z_1\in L_1,\quad Z_2\in\tfrac12L_2,
       \quad Z_3\in\tfrac1{12}L_3.
 \tag{E.81.4}
\]
All coefficients of $Z_4$ have denominator dividing $24$.
Its terms linear in $h$ have denominator dividing $12$;
those linear in $d,e$ have denominator dividing two; and every
nonlinear monomial involves at most one of $d,e$.
These are polynomial identities in the five formal scalar
variables, not estimates obtained by testing finitely many
integral substitutions. The retained program
\path{code/check_class4_denominators.py} expands all expressions
exactly over $\mathbb Q$, keeps the scalar variables formal,
and verifies each assertion. To see that associative coefficients
control the Lie lattice, in degrees two, three and four use the
bases
\[
 [X,Y];\qquad [X,[X,Y]],[Y,[X,Y]];\qquad
 [X,[X,[X,Y]]],[Y,[X,[X,Y]]],[Y,[Y,[X,Y]]].
\]
Their successive leading associative words are
\[
 XY;\qquad XXY,XYY;\qquad XXXY,XXYY,XYYY.
\]
Their coefficients are $\pm1$, and their coefficient matrices
are triangular. Thus these lattices are
saturated in their rational spans. The estimates hold also at
two and three. The universal Baker--Campbell--Hausdorff
coefficients through degree four have respective denominator
bounds $1,2,12,24$, as the same exact expansion verifies.

Write $W_i=(\log V)_i$ and $Z_i=(\log U)_i$.
The ordinary valuation patterns give
\[
 W_1=mX,\quad mY,\quad\text{or }m(X+Y).
\]
The degree-two part of (E.81.2) reads
\[
              (q-1)W_2=-[Z_1,W_1].
 \tag{E.81.5}
\]
In the first pattern the second Kummer class is locally trivial,
so its Frobenius value $\beta$ is $(q-1)u$ for integral $u$.
In the second pattern the first value $\alpha$ is $(q-1)u$.
In the third pattern, the unit ratio tends to $-1$ in the local
Kummer group, so
\[
                 \beta-\alpha=(q-1)(u+\tfrac12),
                 \qquad u\in\mathbb Z_\ell.
\]
This is the actual identity for the Kummer cocycle of $-1$,
$(\chi-1)/2$, with an integral coboundary; it is valid at two.
Consequently
\[
                            W_2\in\tfrac m2L_2.
 \tag{E.81.6}
\]

The path element has $(\log f)_2\in L_2$ and
$(\log f)_3\in\frac12L_3$ by (E.81.3).
Thus its only extra term in the degree-three action on $W_1$
is in $qmL_3$. Taking degree three in
$\exp(\operatorname{ad}Z)\sigma(W)=qW$ now gives
\[
 (q^2-1)W_3\in\tfrac m2L_3.
 \tag{E.81.7}
\]
Indeed the remaining terms, after division by the unit $q$,
are $q[Z_1,W_2]$, $[Z_2,W_1]$ and
$\frac12[Z_1,[Z_1,W_1]]$; (E.81.4) and (E.81.6)
put each in $\frac m2L_3$.
Set
\[
 K=12e_q(q^3-1),\qquad J=e_q(q^3-1).
 \tag{E.81.8}
\]
Signs have no effect on the following ideal inclusions.
Equations (E.81.6)--(E.81.7) show that $W_1,W_2,W_3$ all
lie in $KL_1,KL_2,KL_3$, respectively.

Take the Hall lift $\widetilde V$ to class four with new
weight-four coordinates zero. The coordinates $a,b$ of $V$
are multiples of $K$. The degree-two inverse formula makes
$h$ a multiple of $K$, since its quadratic correction has
denominator two and $12\mid K$. The degree-three inverse
formula gives only
\[
                            d,e\in(K/2)\mathbb Z_\ell.
\]
This possible loss of two is retained. In that formula the
linear $h$ term has denominator two, and all remaining
nonlinear lower-coordinate terms have denominator dividing
$12$, hence are multiples of $K$. By (E.81.3), the degree-four
terms linear in $h,d,e$ lie in $(K/12)L_4$. A nonlinear
term has at least two coordinates, at most one from $d,e$,
and denominator dividing $24$, so it lies in
$(K^2/48)L_4\subseteq(K/12)L_4$. Therefore
\[
                    (\log\widetilde V)_4\in JL_4.
 \tag{E.81.9}
\]

Choose any integral lift $\widetilde U$ of $U$. Its first three
logarithmic components obey (E.81.4); its fourth component
cannot affect conjugation in degree four. The degree-four part
of $\sigma(\log\widetilde V)$ consists of $q^4W_4$,
a term in $(m/2)L_4$ from the degree-three part of $\log f$,
and a term in $(m/2)L_4$ from its degree-two part acting on
$W_2$. All lie in $JL_4$. In applying
$\exp(\operatorname{ad}Z)$, the other possible degree-four
terms are
\[
\begin{split}
 &[Z_1,(\sigma W)_3],\quad[Z_2,(\sigma W)_2],
       \quad[Z_3,(\sigma W)_1],\\
 &\tfrac12[Z_1,[Z_1,(\sigma W)_2]],\quad
   \tfrac12[Z_1,[Z_2,(\sigma W)_1]],\quad
   \tfrac12[Z_2,[Z_1,(\sigma W)_1]],\\
 &\tfrac16[Z_1,[Z_1,[Z_1,(\sigma W)_1]]].
\end{split}
\]
Their denominator losses are at most $2,12,4,6$ as indicated
by (E.81.4). Terms involving $W_1$ contain $m=24e_q(q^2-1)(q^3-1)$;
the others contain $K=12J$. Each term is therefore in $JL_4$.

The two logarithms
$\log(\widetilde U\sigma(\widetilde V)\widetilde U^{-1})$
and $\log(\widetilde V^q)$ agree through degree three by
(E.81.2). Their difference is central, so taking their quotient
introduces no further Baker--Campbell--Hausdorff term. We obtain
\[
 \widetilde U\sigma(\widetilde V)\widetilde U^{-1}
       \widetilde V^{-q}\in J L_4.
 \tag{E.81.10}
\]
In the tame transgression calculation B.22 this central relator
is, up to sign and multiplication by the unit $q$, the residue
of the actual obstruction cocycle. Consequently its unique
normalized residue, divided by $q^3-1$, lies in
$e_q\mathbb Z_\ell^3$. These residues are integral and tend
to zero. At a zero-support prime the section has trivial
inertia, by successive torsion-free tame equations, so its
residue is zero.

At the old primes two and three, the only nonunit factors of
$q^3-1$ are respectively $7$ and $2,13$. The chosen local
class-three point matches make their fourth obstruction classes
zero; thus their normalized residue coordinates are integral.
All other coefficient primes have unit denominators there.
There are finitely many old places, so they cause no convergence
issue. A.22 now places the global class-four obstruction in
$H^2(G_{\mathbb Q,\{\ell,\infty\}},\mathbb Z_\ell(4))^3$.

For odd $\ell$ this last group is zero, by the exact comparison
with $K_6(\mathbb Z)\{\ell\}=0$. For $\ell=2$ its kernel
under real restriction is zero: the comparison in
\citet[VI.9.11, case $n=6$, pp.57--58; VI.10.3, p.61]{CohWeibelKbook}
identifies that kernel with $K_6(\mathbb Z[1/2])\{2\}=0$.
Localizing at $\ell$ does not change the $\ell$-primary part,
since $K_5(\mathbb F_\ell)$ has order $\ell^3-1$.
The same source, VI.8.2, gives the odd-primary comparison.
The real class is zero because E.80 made the class-three
restriction an actual real point section, which has a full
lift. Thus the fixed-ramification obstruction is zero at two
as well. Continuous primitives provide global pro-$\ell$
class-four lifts. Taking their product proves the integral
class-four assertion.
\end{proof}

\begin{result}{E.82}{Proved}{The exact cubic-exponent ordinary recursion, including initial factorial six}
Let $24\mid C$, and let $J\geq6$ be at least every prime
divisor of $C$. With exactly the parameters
\[
 n_j=(j+J)!,\qquad D_j=6Cn_j^3,\qquad
 E_j(Q)=Cn_j(Q^2-1)(Q^3-1),
 \tag{E.82.1}
\]
there are disjoint sequences of distinct primes $q_j,r_j,s_j>3$
for which
\[
 A=3\prod_{j\geq0}\frac{q_j^{E_j(q_j)}}{s_j^{E_j(s_j)}},
 \qquad
 B=-2\prod_{j\geq0}\frac{r_j^{E_j(r_j)}}{s_j^{E_j(s_j)}}
 \tag{E.82.2}
\]
is ordinary at every finite prime and has infinite support.
Its valuation patterns are $(0,1)$ at two, $(1,0)$ at three,
and respectively $(E_j(q_j),0),(0,E_j(r_j))$ and
$(-E_j(s_j),-E_j(s_j))$ at the new primes.
This independently confirms SECTIONS\_OVER\_Q A.62 with its
literal $J=6$ option. With $C=24,J=6$, E.81 gives an actual
integral class-four section above this pair.
\end{result}
\begin{proof}
We verify the small-factorial estimates that are not supplied
by the general $J\geq18$ bound in E.30--E.31, and then give
the resulting recursion with these exact parameters.
Fix a stage and write $n=n_j$, $D=D_j$ and $E=E_j$.
For $p\mid n$ put $v=v_p(n)$, $c_p=v_p(C)$,
$a_2=2$, and $a_p=1$ for odd $p$.
The exact logarithmic identity for $z=1+p^vu$ is
\[
 E(z)\log z=6Cn p^{3v}u^3H(u),
 \quad
 H(u)=\left(1+\frac{p^vu}{2}\right)
       \left(1+p^vu+\frac{p^{2v}u^2}{3}\right)
       \frac{\log(1+p^vu)}{p^vu}.
 \tag{E.82.3}
\]
At two $v\geq4$, and at three $v\geq2$ because $6!\mid n$.
For all other $p\mid n$, $v\geq1$.
The coefficient of $u^{k-1}$ in the final factor has valuation
$(k-1)v-v_p(k)$. At two this is at least $v-1\geq3$ for
$k\geq2$; at three it is at least two; for $p\geq5$
it is at least one. These estimates tend to infinity with $k$.
The two polynomial factors satisfy the same or stronger bounds.
Consequently
\[
 H\in1+p^{1+v_p(3)}\mathbb Z_p\langle u\rangle.
 \tag{E.82.4}
\]
This is the precise coefficient estimate used in E.30.

For clarity, put $u=1+p^{a_p}w$ and $e=v_p(3)$.
After its constant value is removed,
$u^3H(u)$ divided by $p^{a_p+e}$ has the form
$\lambda w+pG(w)$, with $\lambda\in\mathbb Z_p^\times$
and $G$ an integral restricted power series. The binomial
coefficients of $u^3$ prove this directly; (E.82.4) gives
one extra factor $p$ for the nonconstant contribution of
$H-1$. Contraction in $\mathbb Z_p$ proves that this map is
a bijection. Thus the exact shell map is
\[
\begin{split}
 1+p^v+p^{v+a_p}\mathbb Z_p&\xrightarrow{\sim}
 z_0^{E(z_0)}(1+p^{W_p}\mathbb Z_p),\\
 W_p&=c_p+4v+v_p(6)+a_p+v_p(3),\\
 z_0&=1+p^v.
\end{split}
\tag{E.82.5}
\]
Logarithm and exponential are taken in their usual inverse
isometric domains; at two all values here lie in $4\mathbb Z_2$.
The map agrees with the literal integer power on positive integer
inputs, and preserves congruences. Throughout its shell,
\[
 v_p(E(z))=c_p+3v+v_p(6)=v_p(D),
 \qquad W_p\leq2v_p(D).
 \tag{E.82.6}
\]
Indeed the difference in the inequality is
$c_p+2v+v_p(6)-a_p-v_p(3)$, nonnegative also at two and
three with the just proved minimum values of $v$.

If $p\nmid n$, then $p\geq7$ and $p\nmid C$.
On a Teichmuller residue ball with $\omega^2,\omega^3\ne1$
modulo $p$, retain separately a frozen exponent
$e_0\bmod(p-1)$. The map
\[
 z\longmapsto\omega^{e_0}
            \exp(E(z)\log(z/\omega))
 \tag{E.82.7}
\]
is a bijection onto $\omega^{e_0}(1+p\mathbb Z_p)$.
Its logarithm, after $z=\omega(1+pw)$ and division by $p$,
is $E(\omega)w+pG_0(w)$ with unit $E(\omega)$.
This proves the assertion by contraction. It agrees with the
ordinary power when the independently imposed integer congruence
$E(z)\equiv e_0\pmod{p-1}$ holds. At most four nonzero
residues are excluded by $\omega^2=1$ or $\omega^3=1$.
Since $p\geq7$, at least one permitted residue always remains.

Here is the complete induction to which these maps are applied.
Put $B(D)=(D+1)^4$ and $H_j(p)=2v_p(D_j)+j+3$.
Keep disjoint-support coefficients
$a=3u^{24},b=-2v^{24},c=w^{24}$, unit-coordinate local
solutions to $aX^{D_j}+bY^{D_j}=cZ^{D_j}$ at every finite
prime, and $v_p(a+b-c)\geq H_j(p)$ for $p\leq B(D_j)$.
The initial coefficients $(3,-2,1)$ satisfy both conditions
exactly. Choose $P\geq B(D_{j+1})$ including all old support,
and precisions $L_p\geq H_{j+1}(p)$ with $p^{L_p}>P$.
Increase them to retain the full shells in (E.82.5).

Process primes $p\leq P$ increasingly. Earlier congruences
fix each future input modulo $p-1$ and hence its own frozen
exponent $e_i$. The shell valuations (E.82.6), and the exclusion
of square and cube roots of one at primes not dividing $n$, give
\[
 \gcd(e_i,p-1)=\gcd(D,p-1)
       \quad\text{for each of the three input variables}.
 \tag{E.82.8}
\]
All prime factors of $6C$ divide $n$, so there is no omitted
exception to this valuation calculation. Thus all three
finite-field power maps have the same image. Every excluded
residue has order two or three and maps to one.

For $p\mid n$ the three images in (E.82.5) are the same
principal-unit coset. The near-sum invariant and
$H_j(p)\geq W_p$ allow correction of an output with unit
coefficient to solve the exact local equation. Invert its
shell map. For $p\nmid n$, $p\leq B(D)$, choose a common
nonzero power residue different from one if the image has
more than one element. Every preimage is then permitted.
If the power image is trivial, all nonzero residues are
preimages and the four exclusions leave a permitted one.
The common residue and the near-sum invariant allow a
principal-unit correction with a unit coefficient; (E.82.7)
attains the corrected output. For $p>B(D)$ use the maintained
unit-coordinate solution, and scale its three coordinates by
a common unit so that no resulting power value is one modulo
$p$. There are at most three forbidden scaling powers and
$(p-1)/\gcd(D,p-1)>3$. Formula (E.82.8) and (E.82.7)
then give the three exact inputs. This includes old support
primes larger than $B(D)$.

Retain these inputs modulo $p^{L_p}$. Their congruence
preservation gives the required next error precision.
CRT produces three unit classes, each one modulo $n$.
The three fields in (E.31.6) have maximal abelian subfield
$\mathbb Q(\mu_n)$ by its fully proved radical-intersection
argument, which uses only $24\mid n$ and the specified
coefficient forms. They therefore admit simultaneously the
identity on each radical field and its prescribed CRT
cyclotomic automorphism. Chebotarev chooses distinct
$q,r,s>P$ in those classes and respective splitting fields.
Update by $a'=aq^{E(q)},b'=br^{E(r)},c'=cs^{E(s)}$.
Each exponent is a multiple of $24$.

The next local-solubility checks of E.31 use only estimates
already established here. At $p\leq P$, the retained error
is greater than twice the derivative valuation for exponent
$D_{j+1}$, so strong Hensel gives a unit point. At a new prime
$Q$, the exact shell valuations and excluded residue one
at newly entering prime divisors give
$\gcd(D_{j+1},Q-1)=n_j$. The split radical is an
$n_j$th power, as are all new cross-factors, so it is a
$D_{j+1}$th power modulo $Q$; simple-root Hensel solves
using the other two unit coordinates. At the remaining
primes, all coefficients are units and $p>B(D_{j+1})$.
The precise diagonal-curve Weil estimate proved in E.31
gives a nonzero residue point, which again Hensel-lifts.
This proves the induction without enlarging $J$ or changing
any exponent in (E.82.1).

Finally, every update exponent is divisible by $n_j$.
The three global products converge in the power completion.
At each fixed finite prime their valuations eventually
stabilize, their unit components converge, and the increasing
near-sum precision gives nonzero limits $a_p+b_p=c_p$.
Thus $x_p=a_p/c_p$ and $1-x_p=b_p/c_p$ are the two required
ordinary local numbers. Their valuations are exactly the
stated ones. Disjoint new supports give infinite support.
For $C=24$ the coefficients $e_{q_j},e_{r_j},e_{s_j}$ in
(E.81.1) are $n_j$, which tend to zero in every $\mathbb Z_\ell$.
E.81 applies. The valuation at two excludes the cusp-zero
pattern, that at three excludes cusp one, and their unequal
coordinates exclude cusp infinity. Infinite support excludes
a rational-point pair.
\end{proof}

\begin{result}{E.83}{Proved}{The same ordinary ghost has a global class-four section and no locally matching class-two section}
There is an output of the exact recursion E.82 with $C=24,J=6$
satisfying
\[
 (A,B)\bmod8=([3],[-2]),\qquad
                         x_{1129}\equiv828\pmod{1129}.
 \tag{E.83.1}
\]
It admits an actual continuous global integral class-four section,
and its everywhere locally point-or-cusp-matching integral
class-two fibre is empty. In particular its everywhere locally
point-or-cusp-matching class-four fibre is empty as well.
This confirms the full assertions of SECTIONS\_OVER\_Q A.207
and A.209 for the same completed pair. The exclusion is already
seen in the actual geometric order-$32$ quotient of E.70;
its order-$128$ predecessor gives the stated class-four exclusion.
No assertion that the unrestricted class-two or class-four
fibres are empty is made.
\end{result}
\begin{proof}
Modify exactly one local input of the first stage of E.82.
Let $p=1129$, so
\[
 n_0=720,\quad D_0=144\cdot720^3=53747712000,
 \quad p-1=24\cdot47,\quad \gcd(D_0,p-1)=24.
\]
Direct modular exponentiation gives
\[
 59^{D_0}\equiv276,\quad28^{D_0}\equiv978,
 \quad3\cdot276-2\cdot978\equiv1\pmod p.
 \tag{E.83.2}
\]
The integer $D_0$ is a $p$-adic unit. Hensel's lemma therefore
gives the unique $Z_0\in1+p\mathbb Z_p$ satisfying
\[
 Z_0^{D_0}=3\cdot59^{D_0}-2\cdot28^{D_0}.
 \tag{E.83.3}
\]
When the increasing-prime CRT process reaches $p$, the already
fixed congruences determine three possibly different exponent
classes $e_i$ modulo $p-1$. By (E.82.8) each has greatest
common divisor $24$ with $p-1$. Their images are the same
order-$47$ subgroup. Each target $276,978,1$ has exactly
$24$ preimages under each map. At most four nonzero residues
are square or cube roots of one. Thus each target has a
permitted preimage. The frozen-exponent bijections (E.82.7)
attain the exact three outputs
\[
                         59^{D_0},\quad28^{D_0},\quad Z_0^{D_0}.
\]
Their equation is (E.83.3). Retain their inputs to the original
precision $p^{L_p}$. Every exclusion and error precision used
in E.82 is preserved, including those required when later
primes are processed. The radical fields and their CRT
intersections have not changed. The original Chebotarev
construction is therefore valid with this local modification.

The first updated coefficients have residues
$(a_1,b_1,c_1)=(828,302,1)\pmod p$.
At each later stage $p\leq B(D_j)$. If $p\nmid n_j$,
use the common output residue of the small-prime step in
E.82; its principal-unit corrections leave that residue
unchanged. Multiplying the three coefficients by a common
residue preserves $a_j/c_j\pmod p$. Once $p\mid n_j$,
each input is one modulo $p$, so all three multipliers
are one modulo $p$ and again preserve the ratio.
All recursive support primes exceed the processed bound;
none equals $1129$. The ordinary limiting ratio is
therefore $828$ modulo $1129$.

Every update exponent is divisible by $24$, proving the
individual mod-eight classes in (E.83.1). The exact valuations
are those of E.81 with $e_q=n_j$, tending to zero in the
power completion. Its continuous integral class-four lift
exists by E.81. This proves the positive assertion for this
particular pair, rather than for a different fixed-depth output.
Choosing the least permitted residue at unspecified finite
steps, the unique roots on the chosen shells, and then the
least successive primes in the established CRT and radical
classes specifies a complete effective recursion. The
astronomical first prime search has not been executed.

We also check the literal eight-sheet calculation in A.207.
The four quartic branch residues of $C$ at $1129$, in its
original order, are $16,430,699,1113$. The radicand of the
cyclic-eight cover at $t=828$ is $1115$. Substitution in the
actual birational maps (E.70.4), followed by the four quadratic
characters of $x-\theta$, gives
\[
\begin{array}{r|r|r|c}
 z&x&y&\text{four Kummer bits}\\\hline
 35&28&454&0101\\
 44&941&603&1100\\
 235&1101&629&1010\\
 511&188&306&0011\\
 618&188&823&0011\\
 894&1101&500&1010\\
 1085&941&526&1100\\
 1094&28&675&0101
\end{array}
\tag{E.83.4}
\]
The full branch-coordinate Kummer comparison, including the
rational fifth coordinate, was proved in E.70. Here bit zero
means a square unit and bit one a nonsquare unit. Every
entry is a unit, and every root derivative is a unit, so
these data persist on the entire residue disc by simple-root
Hensel lifting. The projection to the quadratic norm pairs
the first and fourth bits, and the second and third bits.
Every row therefore projects to $(1,1)$. E.70 proves that
the entire global projected Selmer image has local vector
$(0,0)$ at this prime. This excludes all eight sheets
without assuming the earlier full two-Selmer computation.
It also proves the originally stated order-$128$ exclusion,
since that arithmetic quotient maps to the order-$32$ one.

For the constructed pair, $AB^3\bmod8=[-24]$.
Any locally point-or-cusp-matching integral class-two section
would consequently map to the empty locus of E.70.
Full local multiplicative Kummer injectivity fixes its
ordinary centre to $x_{1129}$. A cusp cannot give these two
local coordinates: the three cusp patterns would require
$1-x=1$, $x=1$, or $(1-x)/x=-1$, none of which holds.
The map from class four to class two carries local point
and cusp sections to the same type of sections, proving
both exclusions.

The retained program \path{code/check_cubic_recursion.py}
checks primality, all $46\cdot3$ frozen-exponent/target
cases, the initial residues, all eight cover points and
all their full and projected Kummer vectors. Its output
is included. The infinite recursion, continuous lifting,
and all-place projected Selmer bound are provided by the
proofs, not inferred from this finite certificate.
\end{proof}
